\chapter{Anexo C. Documentación Técnica}\label{anexo-c.-documentaciuxf3n-tuxe9cnica}

\textbf{OrbitEngine}: Plataforma SaaS para la Gestión de Procesos Internos en Pymes\\
Versión: 1.0 \textbar{} Abril 2026\\
Universidad Sergio Arboleda, Semillero de Software como Innovación

\begin{center}\rule{0.5\linewidth}{0.5pt}\end{center}

\section{Stack Tecnológico}\label{stack-tecnoluxf3gico}

\subsection{Backend}\label{backend}

{\def\LTcaptype{none} % do not increment counter
\begin{longtable}[]{@{}
  >{\raggedright\arraybackslash}p{(\linewidth - 6\tabcolsep) * \real{0.2283}}
  >{\raggedright\arraybackslash}p{(\linewidth - 6\tabcolsep) * \real{0.1848}}
  >{\raggedright\arraybackslash}p{(\linewidth - 6\tabcolsep) * \real{0.0761}}
  >{\raggedright\arraybackslash}p{(\linewidth - 6\tabcolsep) * \real{0.5109}}@{}}
\toprule\noalign{}
\begin{minipage}[b]{\linewidth}\raggedright
Componente
\end{minipage} & \begin{minipage}[b]{\linewidth}\raggedright
Tecnología
\end{minipage} & \begin{minipage}[b]{\linewidth}\raggedright
Versión
\end{minipage} & \begin{minipage}[b]{\linewidth}\raggedright
Propósito
\end{minipage} \\
\midrule\noalign{}
\endhead
\bottomrule\noalign{}
\endlastfoot
Lenguaje & Python & 3.10+ & Lenguaje principal del backend \\
Framework web & FastAPI & 0.115+ & API REST asíncrona con tipado estricto \\
ORM + Schemas & SQLModel & 0.0.21+ & Modelos compartidos entre SQLAlchemy y Pydantic \\
Motor de BD & SQLAlchemy & 2.x & Capa de abstracción de base de datos \\
Driver PostgreSQL & Psycopg3 & 3.x & Conector PostgreSQL para Python \\
Migraciones & Alembic & 1.x & Gestión de migraciones de esquema \\
Autenticación & PyJWT & 2.x & Generación y verificación de tokens JWT \\
Hash de contraseñas & Pwdlib & 0.2+ & Hashing seguro con bcrypt \\
Validación & Pydantic & 2.x & Schemas de entrada/salida de datos \\
Configuración & Pydantic-Settings & 2.x & Variables de entorno tipadas \\
Linter/Formatter & Ruff & 0.6+ & Análisis estático y formato de código \\
Verificación de tipos & Mypy & 1.x & Chequeo de tipos estático en modo estricto \\
Testing & Pytest + Coverage & 8.x & Pruebas unitarias e integración \\
Monitoring & Sentry (opcional) & N/A & Rastreo de errores en producción \\
\end{longtable}
}

\subsection{Frontend}\label{frontend}

{\def\LTcaptype{none} % do not increment counter
\begin{longtable}[]{@{}
  >{\raggedright\arraybackslash}p{(\linewidth - 6\tabcolsep) * \real{0.1839}}
  >{\raggedright\arraybackslash}p{(\linewidth - 6\tabcolsep) * \real{0.1839}}
  >{\raggedright\arraybackslash}p{(\linewidth - 6\tabcolsep) * \real{0.0805}}
  >{\raggedright\arraybackslash}p{(\linewidth - 6\tabcolsep) * \real{0.5517}}@{}}
\toprule\noalign{}
\begin{minipage}[b]{\linewidth}\raggedright
Componente
\end{minipage} & \begin{minipage}[b]{\linewidth}\raggedright
Tecnología
\end{minipage} & \begin{minipage}[b]{\linewidth}\raggedright
Versión
\end{minipage} & \begin{minipage}[b]{\linewidth}\raggedright
Propósito
\end{minipage} \\
\midrule\noalign{}
\endhead
\bottomrule\noalign{}
\endlastfoot
Lenguaje & TypeScript & 5.x & Tipado estático sobre JavaScript \\
Framework UI & React & 19.x & Librería de componentes reactivos \\
Build tool & Vite & 6.x & Empaquetado y servidor de desarrollo \\
Router & TanStack Router & 1.x & Enrutamiento con tipado por archivo \\
Server state & TanStack Query & 5.x & Cache y sincronización de datos del servidor \\
Tablas & TanStack Table & 8.x & Tablas con sorting/filtering/paginación \\
Formularios & React Hook Form & 7.x & Manejo de formularios con rendimiento optimizado \\
Validación & Zod & 3.x & Schemas de validación con inferencia de tipos \\
UI Components & Shadcn/ui & N/A & Componentes base sobre Radix UI \\
Estilos & Tailwind CSS & 4.x & Utilidades CSS de bajo nivel \\
Gráficos & Recharts & 2.x & Gráficas para el dashboard \\
HTTP Client & Axios (generado) & N/A & Cliente OpenAPI auto-generado \\
Notificaciones & Sonner & 1.x & Toasts y notificaciones \\
Íconos & Lucide React & 0.4+ & Librería de íconos SVG \\
Linter/Formatter & Biome & 1.x & Análisis estático y formato \\
E2E Testing & Playwright & 1.x & Pruebas end-to-end en navegador \\
\end{longtable}
}

\subsection{Infraestructura}\label{infraestructura}

{\def\LTcaptype{none} % do not increment counter
\begin{longtable}[]{@{}
  >{\raggedright\arraybackslash}p{(\linewidth - 4\tabcolsep) * \real{0.1806}}
  >{\raggedright\arraybackslash}p{(\linewidth - 4\tabcolsep) * \real{0.2361}}
  >{\raggedright\arraybackslash}p{(\linewidth - 4\tabcolsep) * \real{0.5833}}@{}}
\toprule\noalign{}
\begin{minipage}[b]{\linewidth}\raggedright
Componente
\end{minipage} & \begin{minipage}[b]{\linewidth}\raggedright
Tecnología
\end{minipage} & \begin{minipage}[b]{\linewidth}\raggedright
Propósito
\end{minipage} \\
\midrule\noalign{}
\endhead
\bottomrule\noalign{}
\endlastfoot
Base de datos & PostgreSQL 18 & Almacenamiento relacional principal \\
Contenedores & Docker Compose v2 & Orquestación de servicios \\
Reverse proxy & Traefik 3.6 & Enrutamiento HTTP/HTTPS + TLS automático \\
CI/CD & GitHub Actions & Pipeline de integración y entrega continua \\
Panel de BD & Adminer & Administración web de PostgreSQL \\
Email (dev) & Mailcatcher & Interceptor SMTP para desarrollo \\
\end{longtable}
}

\section{Estructura del Proyecto}\label{estructura-del-proyecto}

\begin{verbatim}
orbit-engine/
├── .env                        # Variables de entorno (no incluir en git)
├── compose.yml                 # Configuración Docker Compose base
├── compose.override.yml        # Sobreescrituras para desarrollo local
├── compose.traefik.yml         # Traefik para producción
├── pyproject.toml              # Configuración Ruff y Mypy (nivel raíz)
│
├── backend/
│   ├── Dockerfile
│   ├── pyproject.toml          # Dependencias Python (uv)
│   ├── scripts/
│   │   ├── prestart.sh         # Ejecutado por el servicio prestart
│   │   ├── test.sh             # Script de pruebas con cobertura
│   │   ├── lint.sh             # Mypy + Ruff check
│   │   └── format.sh           # Ruff format
│   └── app/
│       ├── main.py             # Punto de entrada FastAPI + registro de routers
│       ├── models.py           # Todos los modelos SQLModel + schemas Pydantic
│       ├── crud.py             # Operaciones de base de datos
│       ├── utils.py            # Helpers (email, generación de tokens)
│       ├── backend_pre_start.py # Script de inicialización
│       ├── initial_data.py     # Seed de datos iniciales (roles, superusuario)
│       ├── api/
│       │   ├── deps.py         # Dependencias FastAPI (auth, sesión, roles)
│       │   └── routes/
│       │       ├── login.py            # Autenticación y recuperación de contraseña
│       │       ├── users.py            # Gestión de usuarios
│       │       ├── organizations.py    # Registro y configuración de organización
│       │       ├── roles.py            # Listado de roles
│       │       ├── categories.py       # CRUD de categorías
│       │       ├── products.py         # CRUD de productos + stock
│       │       ├── customers.py        # CRUD de clientes
│       │       ├── sales.py            # Registro y gestión de ventas
│       │       ├── inventory_movements.py  # Historial de movimientos
│       │       ├── dashboard.py        # Estadísticas del dashboard
│       │       ├── utils.py            # Health check
│       │       └── private.py          # Endpoints de superusuario
│       ├── core/
│       │   ├── config.py       # Settings con Pydantic-Settings
│       │   ├── db.py           # Engine SQLAlchemy
│       │   └── security.py     # JWT, hashing de contraseñas
│       └── alembic/
│           ├── env.py
│           └── versions/
│               ├── 001_initial_schema.py   # organizations, roles, users
│               ├── 002_categories.py
│               ├── 003_products.py
│               ├── 004_customers.py
│               ├── 005_inventory_movements.py
│               └── 006_sales.py
│
├── frontend/
│   ├── Dockerfile
│   ├── package.json
│   ├── vite.config.ts
│   ├── openapi-ts.config.ts    # Configuración del generador de cliente API
│   └── src/
│       ├── main.tsx            # Punto de entrada React
│       ├── routeTree.gen.ts    # Árbol de rutas auto-generado (no editar)
│       ├── utils.ts            # Helpers compartidos
│       ├── client/             # Cliente API auto-generado (no editar)
│       │   ├── index.ts
│       │   ├── sdk.gen.ts      # Funciones de llamada a la API
│       │   ├── types.gen.ts    # Tipos TypeScript desde OpenAPI
│       │   ├── schemas.gen.ts
│       │   └── core/           # Configuración Axios
│       ├── components/
│       │   ├── Admin/          # Componentes de gestión de usuarios
│       │   ├── Common/         # Componentes reutilizables (DataTable, layouts)
│       │   ├── Customers/      # Componentes del módulo de clientes
│       │   ├── Inventory/      # Componentes del módulo de inventario
│       │   ├── Landing/        # Componentes de la página pública
│       │   ├── Pending/        # Componentes de estados de carga
│       │   ├── Sales/          # Componentes del módulo de ventas
│       │   ├── Sidebar/        # Navegación principal
│       │   ├── UserSettings/   # Configuración de perfil y organización
│       │   ├── theme-provider.tsx
│       │   └── ui/             # Componentes Shadcn/ui (no editar)
│       ├── hooks/
│       │   ├── useAuth.ts
│       │   ├── useCopyToClipboard.ts
│       │   ├── useCustomToast.ts
│       │   ├── useDebounce.ts
│       │   ├── useMobile.ts
│       │   └── useScrollAnimation.ts
│       ├── lib/
│       │   └── utils.ts        # cn() helper para Tailwind
│       └── routes/
│           ├── __root.tsx      # Layout raíz con providers
│           ├── index.tsx       # Landing page (/)
│           ├── login.tsx       # Inicio de sesión (/login)
│           ├── signup.tsx      # Registro de usuario (/signup)
│           ├── signup-org.tsx  # Registro de organización (/signup-org)
│           ├── recover-password.tsx
│           ├── reset-password.tsx
│           └── dashboard/
│               ├── index.tsx   # Dashboard KPIs (/dashboard)
│               ├── inventory.tsx
│               ├── sales.tsx
│               ├── customers.tsx
│               ├── admin.tsx
│               └── settings.tsx
│
└── docs/
    └── informe-final/          # Documentación del proyecto de grado
\end{verbatim}

\section{Autenticación y Autorización}\label{autenticaciuxf3n-y-autorizaciuxf3n}

\subsection{Flujo de autenticación}\label{flujo-de-autenticaciuxf3n}

OrbitEngine usa \textbf{JWT Bearer Tokens} siguiendo el estándar OAuth2 Password Flow:

\begin{enumerate}
\def\labelenumi{\arabic{enumi}.}
\tightlist
\item
  El cliente envía credenciales a \texttt{POST\ /api/v1/login/access-token}.
\item
  El backend valida las credenciales, verifica que la cuenta no esté bloqueada.
\item
  Si son válidas, retorna un \texttt{access\_token} JWT.
\item
  El cliente incluye el token en el header \texttt{Authorization:\ Bearer\ \textless{}token\textgreater{}} en cada request.
\item
  El backend decodifica el token, extrae el \texttt{sub} (UUID del usuario) y carga el usuario de la BD.
\item
  La organización del usuario se extrae de \texttt{user.organization\_id}, no del token directamente.
\end{enumerate}

\textbf{Parámetros del token:}

{\def\LTcaptype{none} % do not increment counter
\begin{longtable}[]{@{}ll@{}}
\toprule\noalign{}
Parámetro & Valor \\
\midrule\noalign{}
\endhead
\bottomrule\noalign{}
\endlastfoot
Algoritmo & HS256 \\
Payload \texttt{sub} & UUID del usuario (\texttt{str}) \\
Expiración & 8 días (11,520 minutos) \\
Clave de firma & \texttt{SECRET\_KEY} del entorno \\
\end{longtable}
}

\subsection{Control de acceso basado en roles (RBAC)}\label{control-de-acceso-basado-en-roles-rbac}

Los roles se almacenan en la tabla \texttt{roles} con permisos como array JSONB. La verificación en los endpoints usa la dependencia \texttt{require\_role()}:

\begin{Shaded}
\begin{Highlighting}[]
\CommentTok{\# Solo administradores pueden crear usuarios}
\AttributeTok{@router.post}\NormalTok{(}\StringTok{"/"}\NormalTok{, dependencies}\OperatorTok{=}\NormalTok{[Depends(require\_role(}\StringTok{"admin"}\NormalTok{))])}
\KeywordTok{def}\NormalTok{ create\_user(...) }\OperatorTok{{-}\textgreater{}}\NormalTok{ Any: ...}

\CommentTok{\# Administradores y vendedores pueden registrar ventas}
\AttributeTok{@router.post}\NormalTok{(}\StringTok{"/"}\NormalTok{, dependencies}\OperatorTok{=}\NormalTok{[Depends(require\_role(}\StringTok{"admin"}\NormalTok{, }\StringTok{"seller"}\NormalTok{))])}
\KeywordTok{def}\NormalTok{ create\_sale(...) }\OperatorTok{{-}\textgreater{}}\NormalTok{ Any: ...}
\end{Highlighting}
\end{Shaded}

\textbf{Dependencias de autenticación disponibles:}

{\def\LTcaptype{none} % do not increment counter
\begin{longtable}[]{@{}
  >{\raggedright\arraybackslash}p{(\linewidth - 2\tabcolsep) * \real{0.2877}}
  >{\raggedright\arraybackslash}p{(\linewidth - 2\tabcolsep) * \real{0.7123}}@{}}
\toprule\noalign{}
\begin{minipage}[b]{\linewidth}\raggedright
Tipo anotado
\end{minipage} & \begin{minipage}[b]{\linewidth}\raggedright
Descripción
\end{minipage} \\
\midrule\noalign{}
\endhead
\bottomrule\noalign{}
\endlastfoot
\texttt{CurrentUser} & Usuario autenticado; error 403 si token inválido \\
\texttt{CurrentAdminUser} & Usuario autenticado con rol \texttt{admin}; error 403 si no \\
\texttt{CurrentOrganization} & UUID de organización extraído del \texttt{CurrentUser} \\
\texttt{SessionDep} & Sesión de base de datos SQLModel \\
\end{longtable}
}

\subsection{Protección contra fuerza bruta}\label{protecciuxf3n-contra-fuerza-bruta}

El modelo \texttt{User} incluye campos \texttt{failed\_login\_attempts} y \texttt{locked\_until}. Tras N intentos fallidos consecutivos, la cuenta se bloquea temporalmente.

\section{API REST. Referencia de Endpoints}\label{api-rest.-referencia-de-endpoints}

Todos los endpoints están prefijados con \texttt{/api/v1}. La documentación interactiva se sirve en \texttt{/docs} (Swagger UI) y \texttt{/redoc}.

\subsection{\texorpdfstring{Autenticación (\texttt{/login})}{Autenticación (/login)}}\label{autenticaciuxf3n-login}

{\def\LTcaptype{none} % do not increment counter
\begin{longtable}[]{@{}
  >{\raggedright\arraybackslash}p{(\linewidth - 6\tabcolsep) * \real{0.0857}}
  >{\raggedright\arraybackslash}p{(\linewidth - 6\tabcolsep) * \real{0.4000}}
  >{\raggedright\arraybackslash}p{(\linewidth - 6\tabcolsep) * \real{0.0571}}
  >{\raggedright\arraybackslash}p{(\linewidth - 6\tabcolsep) * \real{0.4571}}@{}}
\toprule\noalign{}
\begin{minipage}[b]{\linewidth}\raggedright
Método
\end{minipage} & \begin{minipage}[b]{\linewidth}\raggedright
Ruta
\end{minipage} & \begin{minipage}[b]{\linewidth}\raggedright
Auth
\end{minipage} & \begin{minipage}[b]{\linewidth}\raggedright
Descripción
\end{minipage} \\
\midrule\noalign{}
\endhead
\bottomrule\noalign{}
\endlastfoot
POST & \texttt{/login/access-token} & No & Obtener token JWT (login) \\
POST & \texttt{/login/test-token} & Sí & Verificar validez del token \\
POST & \texttt{/password-recovery/\{email\}} & No & Solicitar enlace de recuperación \\
POST & \texttt{/reset-password} & No & Restablecer contraseña con token \\
\end{longtable}
}

\subsection{\texorpdfstring{Usuarios (\texttt{/users})}{Usuarios (/users)}}\label{usuarios-users}

{\def\LTcaptype{none} % do not increment counter
\begin{longtable}[]{@{}
  >{\raggedright\arraybackslash}p{(\linewidth - 8\tabcolsep) * \real{0.0789}}
  >{\raggedright\arraybackslash}p{(\linewidth - 8\tabcolsep) * \real{0.2632}}
  >{\raggedright\arraybackslash}p{(\linewidth - 8\tabcolsep) * \real{0.0526}}
  >{\raggedright\arraybackslash}p{(\linewidth - 8\tabcolsep) * \real{0.1316}}
  >{\raggedright\arraybackslash}p{(\linewidth - 8\tabcolsep) * \real{0.4737}}@{}}
\toprule\noalign{}
\begin{minipage}[b]{\linewidth}\raggedright
Método
\end{minipage} & \begin{minipage}[b]{\linewidth}\raggedright
Ruta
\end{minipage} & \begin{minipage}[b]{\linewidth}\raggedright
Auth
\end{minipage} & \begin{minipage}[b]{\linewidth}\raggedright
Rol mínimo
\end{minipage} & \begin{minipage}[b]{\linewidth}\raggedright
Descripción
\end{minipage} \\
\midrule\noalign{}
\endhead
\bottomrule\noalign{}
\endlastfoot
GET & \texttt{/users/} & Sí & admin & Listar usuarios de la organización \\
POST & \texttt{/users/} & Sí & admin & Crear usuario en la organización \\
GET & \texttt{/users/me} & Sí & Cualquiera & Obtener perfil propio \\
PATCH & \texttt{/users/me} & Sí & Cualquiera & Actualizar nombre/apellido/teléfono \\
DELETE & \texttt{/users/me} & Sí & Cualquiera & Eliminar cuenta propia (soft delete) \\
PATCH & \texttt{/users/me/password} & Sí & Cualquiera & Cambiar contraseña \\
GET & \texttt{/users/\{user\_id\}} & Sí & admin & Obtener usuario por ID \\
PATCH & \texttt{/users/\{user\_id\}} & Sí & admin & Actualizar usuario \\
DELETE & \texttt{/users/\{user\_id\}} & Sí & admin & Eliminar usuario \\
\end{longtable}
}

\subsection{\texorpdfstring{Organizaciones (\texttt{/organizations})}{Organizaciones (/organizations)}}\label{organizaciones-organizations}

{\def\LTcaptype{none} % do not increment counter
\begin{longtable}[]{@{}
  >{\raggedright\arraybackslash}p{(\linewidth - 8\tabcolsep) * \real{0.0952}}
  >{\raggedright\arraybackslash}p{(\linewidth - 8\tabcolsep) * \real{0.2667}}
  >{\raggedright\arraybackslash}p{(\linewidth - 8\tabcolsep) * \real{0.0571}}
  >{\raggedright\arraybackslash}p{(\linewidth - 8\tabcolsep) * \real{0.1238}}
  >{\raggedright\arraybackslash}p{(\linewidth - 8\tabcolsep) * \real{0.4286}}@{}}
\toprule\noalign{}
\multicolumn{5}{@{}>{\raggedright\arraybackslash}p{(\linewidth - 8\tabcolsep) * \real{0.9714} + 8\tabcolsep}@{}}{%
\begin{minipage}[b]{\linewidth}\raggedright
Método \textbar{} Ruta \textbar{} Auth\textbar{} Rol mínimo \textbar{} Descripción \textbar{}
\end{minipage}} \\
\midrule\noalign{}
\endhead
\bottomrule\noalign{}
\endlastfoot
\multicolumn{3}{@{}>{\raggedright\arraybackslash}p{(\linewidth - 8\tabcolsep) * \real{0.4190} + 4\tabcolsep}}{%
GET \textbar{} /organizations/me \textbar{} Sí \textbar{} Cualquiera} & \multicolumn{2}{>{\raggedright\arraybackslash}p{(\linewidth - 8\tabcolsep) * \real{0.5524} + 2\tabcolsep}@{}}{%
Ver datos de la organización actual \textbar{}} \\
\multicolumn{5}{@{}>{\raggedright\arraybackslash}p{(\linewidth - 8\tabcolsep) * \real{0.9714} + 8\tabcolsep}@{}}{%
PATCH \textbar{} /organizations/me \textbar{} Sí \textbar{} admin \textbar{} Actualizar nombre/slug/descripción \textbar{}} \\
\multicolumn{5}{@{}>{\raggedright\arraybackslash}p{(\linewidth - 8\tabcolsep) * \real{0.9714} + 8\tabcolsep}@{}}{%
POST \textbar{} /organizations/signup \textbar{} No \textbar{} N/A \textbar{} Registro de nueva organización + admin \textbar{}} \\
\end{longtable}
}

\subsection{\texorpdfstring{Roles (\texttt{/roles})}{Roles (/roles)}}\label{roles-roles}

{\def\LTcaptype{none} % do not increment counter
\begin{longtable}[]{@{}
  >{\raggedright\arraybackslash}p{(\linewidth - 6\tabcolsep) * \real{0.0968}}
  >{\raggedright\arraybackslash}p{(\linewidth - 6\tabcolsep) * \real{0.1452}}
  >{\raggedright\arraybackslash}p{(\linewidth - 6\tabcolsep) * \real{0.0645}}
  >{\raggedright\arraybackslash}p{(\linewidth - 6\tabcolsep) * \real{0.6935}}@{}}
\toprule\noalign{}
\begin{minipage}[b]{\linewidth}\raggedright
Método
\end{minipage} & \begin{minipage}[b]{\linewidth}\raggedright
Ruta
\end{minipage} & \begin{minipage}[b]{\linewidth}\raggedright
Auth
\end{minipage} & \begin{minipage}[b]{\linewidth}\raggedright
Descripción
\end{minipage} \\
\midrule\noalign{}
\endhead
\bottomrule\noalign{}
\endlastfoot
GET & \texttt{/roles/} & Sí & Listar roles disponibles en la organización \\
\end{longtable}
}

\subsection{\texorpdfstring{Categorías (\texttt{/categories})}{Categorías (/categories)}}\label{categoruxedas-categories}

{\def\LTcaptype{none} % do not increment counter
\begin{longtable}[]{@{}
  >{\raggedright\arraybackslash}p{(\linewidth - 8\tabcolsep) * \real{0.0690}}
  >{\raggedright\arraybackslash}p{(\linewidth - 8\tabcolsep) * \real{0.2069}}
  >{\raggedright\arraybackslash}p{(\linewidth - 8\tabcolsep) * \real{0.0460}}
  >{\raggedright\arraybackslash}p{(\linewidth - 8\tabcolsep) * \real{0.1149}}
  >{\raggedright\arraybackslash}p{(\linewidth - 8\tabcolsep) * \real{0.5632}}@{}}
\toprule\noalign{}
\begin{minipage}[b]{\linewidth}\raggedright
Método
\end{minipage} & \begin{minipage}[b]{\linewidth}\raggedright
Ruta
\end{minipage} & \begin{minipage}[b]{\linewidth}\raggedright
Auth
\end{minipage} & \begin{minipage}[b]{\linewidth}\raggedright
Rol mínimo
\end{minipage} & \begin{minipage}[b]{\linewidth}\raggedright
Descripción
\end{minipage} \\
\midrule\noalign{}
\endhead
\bottomrule\noalign{}
\endlastfoot
GET & \texttt{/categories/} & Sí & Cualquiera & Listar categorías (soporta \texttt{search}, \texttt{is\_active}) \\
POST & \texttt{/categories/} & Sí & admin & Crear categoría \\
GET & \texttt{/categories/\{id\}} & Sí & Cualquiera & Obtener categoría por ID \\
PATCH & \texttt{/categories/\{id\}} & Sí & admin & Actualizar categoría \\
DELETE & \texttt{/categories/\{id\}} & Sí & admin & Eliminar / desactivar categoría \\
\end{longtable}
}

\subsection{\texorpdfstring{Productos (\texttt{/products})}{Productos (/products)}}\label{productos-products}

{\def\LTcaptype{none} % do not increment counter
\begin{longtable}[]{@{}
  >{\raggedright\arraybackslash}p{(\linewidth - 8\tabcolsep) * \real{0.0714}}
  >{\raggedright\arraybackslash}p{(\linewidth - 8\tabcolsep) * \real{0.2540}}
  >{\raggedright\arraybackslash}p{(\linewidth - 8\tabcolsep) * \real{0.0476}}
  >{\raggedright\arraybackslash}p{(\linewidth - 8\tabcolsep) * \real{0.1032}}
  >{\raggedright\arraybackslash}p{(\linewidth - 8\tabcolsep) * \real{0.5000}}@{}}
\toprule\noalign{}
\begin{minipage}[b]{\linewidth}\raggedright
Método
\end{minipage} & \multicolumn{4}{>{\raggedright\arraybackslash}p{(\linewidth - 8\tabcolsep) * \real{0.9048} + 6\tabcolsep}@{}}{%
\begin{minipage}[b]{\linewidth}\raggedright
Ruta \textbar{} Auth\textbar{} Rol mínimo \textbar{} Descripción \textbar{}
\end{minipage}} \\
\midrule\noalign{}
\endhead
\bottomrule\noalign{}
\endlastfoot
\multicolumn{5}{@{}>{\raggedright\arraybackslash}p{(\linewidth - 8\tabcolsep) * \real{0.9762} + 8\tabcolsep}@{}}{%
GET \textbar{} /products/ \textbar{} Sí \textbar{} Cualquiera \textbar{} Listar productos (search, is\_active, category\_id, sort\_by) \textbar{}} \\
\multicolumn{5}{@{}>{\raggedright\arraybackslash}p{(\linewidth - 8\tabcolsep) * \real{0.9762} + 8\tabcolsep}@{}}{%
POST \textbar{} /products/ \textbar{} Sí \textbar{} admin \textbar{} Crear producto \textbar{}} \\
\multicolumn{2}{@{}>{\raggedright\arraybackslash}p{(\linewidth - 8\tabcolsep) * \real{0.3254} + 2\tabcolsep}}{%
GET \textbar{} /products/\{id\} \textbar{} Sí \textbar{} Cualquiera} & \multicolumn{3}{>{\raggedright\arraybackslash}p{(\linewidth - 8\tabcolsep) * \real{0.6508} + 4\tabcolsep}@{}}{%
Obtener producto por ID \textbar{}} \\
\multicolumn{4}{@{}>{\raggedright\arraybackslash}p{(\linewidth - 8\tabcolsep) * \real{0.4762} + 6\tabcolsep}}{%
PATCH \textbar{} /products/\{id\} \textbar{} Sí \textbar{} admin \textbar{} Actualizar producto} & \\
DELETE & \multicolumn{4}{>{\raggedright\arraybackslash}p{(\linewidth - 8\tabcolsep) * \real{0.9048} + 6\tabcolsep}@{}}{%
/products/\{id\} \textbar{} Sí \textbar{} admin \textbar{} Eliminar / desactivar producto \textbar{}} \\
\multicolumn{5}{@{}>{\raggedright\arraybackslash}p{(\linewidth - 8\tabcolsep) * \real{0.9762} + 8\tabcolsep}@{}}{%
POST \textbar{} /products/\{id\}/stock-adjust \textbar{} Sí \textbar{} admin \textbar{} Ajuste manual de stock \textbar{}} \\
\multicolumn{5}{@{}>{\raggedright\arraybackslash}p{(\linewidth - 8\tabcolsep) * \real{0.9762} + 8\tabcolsep}@{}}{%
GET \textbar{} /products/\{id\}/movements \textbar{} Sí \textbar{} Cualquiera \textbar{} Historial de movimientos del producto \textbar{}} \\
\end{longtable}
}

\subsection{\texorpdfstring{Clientes (\texttt{/customers})}{Clientes (/customers)}}\label{clientes-customers}

{\def\LTcaptype{none} % do not increment counter
\begin{longtable}[]{@{}
  >{\raggedright\arraybackslash}p{(\linewidth - 8\tabcolsep) * \real{0.0732}}
  >{\raggedright\arraybackslash}p{(\linewidth - 8\tabcolsep) * \real{0.2805}}
  >{\raggedright\arraybackslash}p{(\linewidth - 8\tabcolsep) * \real{0.0488}}
  >{\raggedright\arraybackslash}p{(\linewidth - 8\tabcolsep) * \real{0.1220}}
  >{\raggedright\arraybackslash}p{(\linewidth - 8\tabcolsep) * \real{0.4756}}@{}}
\toprule\noalign{}
\begin{minipage}[b]{\linewidth}\raggedright
Método
\end{minipage} & \begin{minipage}[b]{\linewidth}\raggedright
Ruta
\end{minipage} & \begin{minipage}[b]{\linewidth}\raggedright
Auth
\end{minipage} & \begin{minipage}[b]{\linewidth}\raggedright
Rol mínimo
\end{minipage} & \begin{minipage}[b]{\linewidth}\raggedright
Descripción
\end{minipage} \\
\midrule\noalign{}
\endhead
\bottomrule\noalign{}
\endlastfoot
GET & \texttt{/customers/} & Sí & Cualquiera & Listar clientes (\texttt{search}, \texttt{is\_active}) \\
POST & \texttt{/customers/} & Sí & Cualquiera & Registrar cliente \\
GET & \texttt{/customers/\{id\}} & Sí & Cualquiera & Obtener cliente por ID \\
PATCH & \texttt{/customers/\{id\}} & Sí & Cualquiera & Actualizar cliente \\
DELETE & \texttt{/customers/\{id\}} & Sí & admin & Desactivar cliente \\
GET & \texttt{/customers/\{id\}/sales} & Sí & Cualquiera & Historial de compras del cliente \\
\end{longtable}
}

\subsection{\texorpdfstring{Ventas (\texttt{/sales})}{Ventas (/sales)}}\label{ventas-sales}

{\def\LTcaptype{none} % do not increment counter
\begin{longtable}[]{@{}
  >{\raggedright\arraybackslash}p{(\linewidth - 8\tabcolsep) * \real{0.0583}}
  >{\raggedright\arraybackslash}p{(\linewidth - 8\tabcolsep) * \real{0.1942}}
  >{\raggedright\arraybackslash}p{(\linewidth - 8\tabcolsep) * \real{0.0388}}
  >{\raggedright\arraybackslash}p{(\linewidth - 8\tabcolsep) * \real{0.0971}}
  >{\raggedright\arraybackslash}p{(\linewidth - 8\tabcolsep) * \real{0.6117}}@{}}
\toprule\noalign{}
\begin{minipage}[b]{\linewidth}\raggedright
Método
\end{minipage} & \begin{minipage}[b]{\linewidth}\raggedright
Ruta
\end{minipage} & \begin{minipage}[b]{\linewidth}\raggedright
Auth
\end{minipage} & \begin{minipage}[b]{\linewidth}\raggedright
Rol mínimo
\end{minipage} & \begin{minipage}[b]{\linewidth}\raggedright
Descripción
\end{minipage} \\
\midrule\noalign{}
\endhead
\bottomrule\noalign{}
\endlastfoot
GET & \texttt{/sales/} & Sí & Cualquiera & Listar ventas (\texttt{search}, \texttt{status}, \texttt{payment\_method}, \texttt{sort\_by}) \\
POST & \texttt{/sales/} & Sí & Cualquiera & Registrar venta (descuenta stock automáticamente) \\
GET & \texttt{/sales/\{id\}} & Sí & Cualquiera & Obtener venta por ID \\
POST & \texttt{/sales/\{id\}/cancel} & Sí & admin & Cancelar venta (revierte stock) \\
\end{longtable}
}

\subsection{\texorpdfstring{Movimientos de inventario (\texttt{/inventory-movements})}{Movimientos de inventario (/inventory-movements)}}\label{movimientos-de-inventario-inventory-movements}

{\def\LTcaptype{none} % do not increment counter
\begin{longtable}[]{@{}
  >{\raggedright\arraybackslash}p{(\linewidth - 6\tabcolsep) * \real{0.0638}}
  >{\raggedright\arraybackslash}p{(\linewidth - 6\tabcolsep) * \real{0.2447}}
  >{\raggedright\arraybackslash}p{(\linewidth - 6\tabcolsep) * \real{0.0426}}
  >{\raggedright\arraybackslash}p{(\linewidth - 6\tabcolsep) * \real{0.6489}}@{}}
\toprule\noalign{}
\begin{minipage}[b]{\linewidth}\raggedright
Método
\end{minipage} & \begin{minipage}[b]{\linewidth}\raggedright
Ruta
\end{minipage} & \begin{minipage}[b]{\linewidth}\raggedright
Auth
\end{minipage} & \begin{minipage}[b]{\linewidth}\raggedright
Descripción
\end{minipage} \\
\midrule\noalign{}
\endhead
\bottomrule\noalign{}
\endlastfoot
GET & \texttt{/inventory-movements/} & Sí & Listar movimientos (\texttt{product\_id}, \texttt{movement\_type}, \texttt{sort\_by}) \\
\end{longtable}
}

\subsection{\texorpdfstring{Dashboard (\texttt{/dashboard})}{Dashboard (/dashboard)}}\label{dashboard-dashboard}

{\def\LTcaptype{none} % do not increment counter
\begin{longtable}[]{@{}
  >{\raggedright\arraybackslash}p{(\linewidth - 6\tabcolsep) * \real{0.0833}}
  >{\raggedright\arraybackslash}p{(\linewidth - 6\tabcolsep) * \real{0.2500}}
  >{\raggedright\arraybackslash}p{(\linewidth - 6\tabcolsep) * \real{0.0556}}
  >{\raggedright\arraybackslash}p{(\linewidth - 6\tabcolsep) * \real{0.6111}}@{}}
\toprule\noalign{}
\begin{minipage}[b]{\linewidth}\raggedright
Método
\end{minipage} & \begin{minipage}[b]{\linewidth}\raggedright
Ruta
\end{minipage} & \begin{minipage}[b]{\linewidth}\raggedright
Auth
\end{minipage} & \begin{minipage}[b]{\linewidth}\raggedright
Descripción
\end{minipage} \\
\midrule\noalign{}
\endhead
\bottomrule\noalign{}
\endlastfoot
GET & \texttt{/dashboard/stats} & Sí & Estadísticas consolidadas de la organización \\
\end{longtable}
}

\textbf{Respuesta de \texttt{/dashboard/stats}:}

\begin{Shaded}
\begin{Highlighting}[]
\FunctionTok{\{}
  \DataTypeTok{"sales\_today"}\FunctionTok{:} \FunctionTok{\{} \DataTypeTok{"count"}\FunctionTok{:} \DecValTok{5}\FunctionTok{,} \DataTypeTok{"total"}\FunctionTok{:} \StringTok{"250000.00"} \FunctionTok{\},}
  \DataTypeTok{"sales\_month"}\FunctionTok{:} \FunctionTok{\{} \DataTypeTok{"count"}\FunctionTok{:} \DecValTok{42}\FunctionTok{,} \DataTypeTok{"total"}\FunctionTok{:} \StringTok{"1850000.00"} \FunctionTok{\},}
  \DataTypeTok{"low\_stock\_count"}\FunctionTok{:} \DecValTok{3}\FunctionTok{,}
  \DataTypeTok{"average\_ticket"}\FunctionTok{:} \StringTok{"44047.62"}\FunctionTok{,}
  \DataTypeTok{"top\_products"}\FunctionTok{:} \OtherTok{[}
    \FunctionTok{\{}
      \DataTypeTok{"product\_id"}\FunctionTok{:} \StringTok{"..."}\FunctionTok{,}
      \DataTypeTok{"product\_name"}\FunctionTok{:} \StringTok{"Producto A"}\FunctionTok{,}
      \DataTypeTok{"quantity\_sold"}\FunctionTok{:} \DecValTok{15}\FunctionTok{,}
      \DataTypeTok{"revenue"}\FunctionTok{:} \StringTok{"450000.00"}
    \FunctionTok{\}}
  \OtherTok{]}\FunctionTok{,}
  \DataTypeTok{"sales\_by\_day"}\FunctionTok{:} \OtherTok{[}\FunctionTok{\{} \DataTypeTok{"date"}\FunctionTok{:} \StringTok{"2026{-}04{-}01"}\FunctionTok{,} \DataTypeTok{"count"}\FunctionTok{:} \DecValTok{3}\FunctionTok{,} \DataTypeTok{"total"}\FunctionTok{:} \StringTok{"120000.00"} \FunctionTok{\}}\OtherTok{]}
\FunctionTok{\}}
\end{Highlighting}
\end{Shaded}

\subsection{Utilidades}\label{utilidades}

{\def\LTcaptype{none} % do not increment counter
\begin{longtable}[]{@{}
  >{\raggedright\arraybackslash}p{(\linewidth - 6\tabcolsep) * \real{0.0923}}
  >{\raggedright\arraybackslash}p{(\linewidth - 6\tabcolsep) * \real{0.3385}}
  >{\raggedright\arraybackslash}p{(\linewidth - 6\tabcolsep) * \real{0.0615}}
  >{\raggedright\arraybackslash}p{(\linewidth - 6\tabcolsep) * \real{0.5077}}@{}}
\toprule\noalign{}
\begin{minipage}[b]{\linewidth}\raggedright
Método
\end{minipage} & \begin{minipage}[b]{\linewidth}\raggedright
Ruta
\end{minipage} & \begin{minipage}[b]{\linewidth}\raggedright
Auth
\end{minipage} & \begin{minipage}[b]{\linewidth}\raggedright
Descripción
\end{minipage} \\
\midrule\noalign{}
\endhead
\bottomrule\noalign{}
\endlastfoot
GET & \texttt{/utils/health-check/} & No & Health check del servicio backend \\
\end{longtable}
}

\section{Modelos de Datos}\label{modelos-de-datos}

\subsection{Organization}\label{organization}

\begin{verbatim}
organizations
├── id              UUID           PK
├── name            VARCHAR(255)   NOT NULL
├── slug            VARCHAR(100)   UNIQUE, INDEX
├── description     TEXT           NULLABLE
├── logo_url        VARCHAR(500)   NULLABLE
├── is_active       BOOLEAN        DEFAULT true
├── created_at      TIMESTAMPTZ    NOT NULL
└── updated_at      TIMESTAMPTZ    NOT NULL

Índices: idx_organizations_slug (unique), idx_organizations_is_active
\end{verbatim}

\textbf{Validación del slug:} expresión regular \texttt{\^{}{[}a-z0-9-{]}\{3,50\}\$} (solo minúsculas, números y guiones).

\subsection{Role}\label{role}

\begin{verbatim}
roles
├── id              INTEGER        PK (autoincrement)
├── name            VARCHAR(50)    UNIQUE, INDEX
├── description     TEXT           NULLABLE
├── permissions     JSONB          DEFAULT []
└── created_at      TIMESTAMPTZ    NOT NULL

Índices: idx_roles_name (unique)
\end{verbatim}

\subsection{User}\label{user}

\begin{verbatim}
users
├── id                     UUID           PK
├── organization_id        UUID           FK → organizations.id, INDEX
├── role_id                INTEGER        FK → roles.id, INDEX
├── email                  VARCHAR(255)   INDEX
├── first_name             VARCHAR(100)   NOT NULL
├── last_name              VARCHAR(100)   NOT NULL
├── phone                  VARCHAR(50)    NULLABLE
├── avatar_url             VARCHAR(500)   NULLABLE
├── hashed_password        VARCHAR(255)   NOT NULL
├── is_active              BOOLEAN        DEFAULT true
├── is_verified            BOOLEAN        DEFAULT false
├── last_login_at          TIMESTAMPTZ    NULLABLE
├── failed_login_attempts  INTEGER        DEFAULT 0
├── locked_until           TIMESTAMPTZ    NULLABLE
├── created_at             TIMESTAMPTZ    NOT NULL
├── updated_at             TIMESTAMPTZ    NOT NULL
└── deleted_at             TIMESTAMPTZ    NULLABLE (soft delete)

Restricciones: UNIQUE(organization_id, email)
Índices: idx_users_organization_id, idx_users_role_id, idx_users_is_active
\end{verbatim}

\subsection{Category}\label{category}

\begin{verbatim}
categories
├── id              UUID           PK
├── organization_id UUID           FK → organizations.id, INDEX
├── parent_id       UUID           FK → categories.id, NULLABLE (jerarquía)
├── name            VARCHAR(255)   NOT NULL
├── description     TEXT           NULLABLE
├── is_active       BOOLEAN        DEFAULT true
├── created_at      TIMESTAMPTZ    NOT NULL
├── updated_at      TIMESTAMPTZ    NOT NULL
└── deleted_at      TIMESTAMPTZ    NULLABLE (soft delete)

Restricciones: UNIQUE(organization_id, name, parent_id)
Índices: idx_categories_organization_id, idx_categories_parent_id, idx_categories_is_active
\end{verbatim}

\subsection{Product}\label{product}

\begin{verbatim}
products
├── id              UUID               PK
├── organization_id UUID               FK → organizations.id, INDEX
├── category_id     UUID               FK → categories.id, NULLABLE
├── name            VARCHAR(255)       NOT NULL
├── sku             VARCHAR(100)       NOT NULL
├── description     TEXT               NULLABLE
├── image_url       VARCHAR(500)       NULLABLE
├── cost_price      NUMERIC(12,2)      DEFAULT 0
├── sale_price      NUMERIC(12,2)      DEFAULT 0
├── stock_quantity  INTEGER            DEFAULT 0
├── stock_min       INTEGER            DEFAULT 0
├── stock_max       INTEGER            NULLABLE
├── unit            VARCHAR(50)        DEFAULT 'unit'
├── barcode         VARCHAR(100)       NULLABLE
├── is_active       BOOLEAN            DEFAULT true
├── created_at      TIMESTAMPTZ        NOT NULL
├── updated_at      TIMESTAMPTZ        NOT NULL
└── deleted_at      TIMESTAMPTZ        NULLABLE (soft delete)

Restricciones: UNIQUE(organization_id, sku)
Índices: idx_products_organization_id, idx_products_category_id, idx_products_is_active, idx_products_sku
\end{verbatim}

\textbf{Stock bajo:} un producto está en alerta cuando \texttt{stock\_quantity\ ≤\ stock\_min}.

\subsection{Customer}\label{customer}

\begin{verbatim}
customers
├── id                UUID           PK
├── organization_id   UUID           FK → organizations.id, INDEX
├── document_type     VARCHAR(50)    NOT NULL (CC, NIT, Pasaporte, Otro)
├── document_number   VARCHAR(50)    NOT NULL
├── first_name        VARCHAR(100)   NOT NULL
├── last_name         VARCHAR(100)   NOT NULL
├── email             VARCHAR(255)   NULLABLE, INDEX
├── phone             VARCHAR(50)    NULLABLE, INDEX
├── address           TEXT           NULLABLE
├── city              VARCHAR(100)   NULLABLE
├── country           VARCHAR(100)   NULLABLE
├── notes             TEXT           NULLABLE
├── is_active         BOOLEAN        DEFAULT true
├── total_purchases   NUMERIC(12,2)  DEFAULT 0 (actualizado automáticamente)
├── purchases_count   INTEGER        DEFAULT 0 (actualizado automáticamente)
├── last_purchase_at  TIMESTAMPTZ    NULLABLE (actualizado automáticamente)
├── created_at        TIMESTAMPTZ    NOT NULL
├── updated_at        TIMESTAMPTZ    NOT NULL
└── deleted_at        TIMESTAMPTZ    NULLABLE (soft delete)

Restricciones: UNIQUE(organization_id, document_number)
Índices: idx_customers_organization_id, idx_customers_email, idx_customers_phone, idx_customers_is_active
\end{verbatim}

\subsection{Sale}\label{sale}

\begin{verbatim}
sales
├── id                   UUID           PK
├── organization_id      UUID           FK → organizations.id, INDEX
├── customer_id          UUID           FK → customers.id, NULLABLE
├── user_id              UUID           FK → users.id, INDEX
├── invoice_number       VARCHAR(100)   NOT NULL
├── sale_date            TIMESTAMPTZ    NOT NULL
├── subtotal             NUMERIC(12,2)  DEFAULT 0
├── discount             NUMERIC(12,2)  DEFAULT 0
├── tax                  NUMERIC(12,2)  DEFAULT 0
├── total                NUMERIC(12,2)  DEFAULT 0
├── payment_method       VARCHAR(50)    DEFAULT 'cash' (cash/card/transfer/other)
├── status               VARCHAR(50)    DEFAULT 'completed' (completed/cancelled/pending)
├── notes                TEXT           NULLABLE
├── cancelled_at         TIMESTAMPTZ    NULLABLE
├── cancelled_by         UUID           FK → users.id, NULLABLE
├── cancellation_reason  TEXT           NULLABLE
├── created_at           TIMESTAMPTZ    NOT NULL
└── updated_at           TIMESTAMPTZ    NOT NULL

Restricciones: UNIQUE(organization_id, invoice_number)
Índices: idx_sales_organization_id, idx_sales_customer_id, idx_sales_user_id,
         idx_sales_status, idx_sales_sale_date, idx_sales_invoice_number
\end{verbatim}

\textbf{Fórmula:} \texttt{total\ =\ subtotal\ -\ discount\ +\ tax}

\subsection{SaleItem}\label{saleitem}

\begin{verbatim}
sale_items
├── id            UUID           PK
├── sale_id       UUID           FK → sales.id, INDEX
├── product_id    UUID           FK → products.id, INDEX
├── product_name  VARCHAR(255)   NOT NULL (snapshot al momento de la venta)
├── product_sku   VARCHAR(100)   NOT NULL (snapshot)
├── quantity      INTEGER        NOT NULL
├── unit_price    NUMERIC(12,2)  NOT NULL
├── subtotal      NUMERIC(12,2)  NOT NULL
└── created_at    TIMESTAMPTZ    NOT NULL

Índices: idx_sale_items_sale_id, idx_sale_items_product_id
\end{verbatim}

\begin{quote}
Los campos \texttt{product\_name} y \texttt{product\_sku} son snapshots del producto en el momento de la venta. Si el producto se edita posteriormente, el historial de ventas conserva los datos originales.
\end{quote}

\subsection{InventoryMovement}\label{inventorymovement}

\begin{verbatim}
inventory_movements
├── id              UUID           PK
├── organization_id UUID           FK → organizations.id, INDEX
├── product_id      UUID           FK → products.id, INDEX
├── user_id         UUID           FK → users.id, INDEX
├── movement_type   VARCHAR(50)    NOT NULL (sale/adjustment/return/purchase)
├── quantity        INTEGER        NOT NULL (positivo = entrada, negativo = salida)
├── previous_stock  INTEGER        NOT NULL
├── new_stock       INTEGER        NOT NULL
├── reference_id    UUID           NULLABLE (FK al objeto de origen)
├── reference_type  VARCHAR(50)    NULLABLE (sale/adjustment/return)
├── reason          TEXT           NULLABLE
└── created_at      TIMESTAMPTZ    NOT NULL

Índices: idx_inventory_movements_organization_id, idx_inventory_movements_product_id,
         idx_inventory_movements_user_id, idx_inventory_movements_movement_type,
         idx_inventory_movements_created_at,
         idx_inventory_movements_reference (reference_type, reference_id)
\end{verbatim}

\section{Variables de Entorno}\label{variables-de-entorno}

La configuración se carga desde el archivo \texttt{.env} usando Pydantic-Settings. Todas las variables se validan al iniciar el servicio.

\subsection{Variables requeridas}\label{variables-requeridas}

{\def\LTcaptype{none} % do not increment counter
\begin{longtable}[]{@{}
  >{\raggedright\arraybackslash}p{(\linewidth - 4\tabcolsep) * \real{0.3662}}
  >{\raggedright\arraybackslash}p{(\linewidth - 4\tabcolsep) * \real{0.1408}}
  >{\raggedright\arraybackslash}p{(\linewidth - 4\tabcolsep) * \real{0.4930}}@{}}
\toprule\noalign{}
\begin{minipage}[b]{\linewidth}\raggedright
Variable
\end{minipage} & \begin{minipage}[b]{\linewidth}\raggedright
Tipo
\end{minipage} & \begin{minipage}[b]{\linewidth}\raggedright
Descripción
\end{minipage} \\
\midrule\noalign{}
\endhead
\bottomrule\noalign{}
\endlastfoot
\texttt{PROJECT\_NAME} & \texttt{str} & Nombre del proyecto \\
\texttt{SECRET\_KEY} & \texttt{str} & Clave para firmar los JWT \\
\texttt{POSTGRES\_SERVER} & \texttt{str} & Host del servidor PostgreSQL \\
\texttt{POSTGRES\_USER} & \texttt{str} & Usuario de PostgreSQL \\
\texttt{POSTGRES\_DB} & \texttt{str} & Nombre de la base de datos \\
\texttt{FIRST\_SUPERUSER} & \texttt{EmailStr} & Email del superusuario inicial \\
\texttt{FIRST\_SUPERUSER\_PASSWORD} & \texttt{str} & Contraseña del superusuario inicial \\
\end{longtable}
}

\subsection{Variables opcionales con valores por defecto}\label{variables-opcionales-con-valores-por-defecto}

{\def\LTcaptype{none} % do not increment counter
\begin{longtable}[]{@{}
  >{\raggedright\arraybackslash}p{(\linewidth - 6\tabcolsep) * \real{0.2353}}
  >{\raggedright\arraybackslash}p{(\linewidth - 6\tabcolsep) * \real{0.1765}}
  >{\raggedright\arraybackslash}p{(\linewidth - 6\tabcolsep) * \real{0.2206}}
  >{\raggedright\arraybackslash}p{(\linewidth - 6\tabcolsep) * \real{0.3529}}@{}}
\toprule\noalign{}
\multicolumn{4}{@{}>{\raggedright\arraybackslash}p{(\linewidth - 6\tabcolsep) * \real{0.9853} + 6\tabcolsep}@{}}{%
\begin{minipage}[b]{\linewidth}\raggedright
Variable \textbar{} Tipo \textbar{} Default \textbar{} Descripción \textbar{}
\end{minipage}} \\
\midrule\noalign{}
\endhead
\bottomrule\noalign{}
\endlastfoot
\multicolumn{4}{@{}>{\raggedright\arraybackslash}p{(\linewidth - 6\tabcolsep) * \real{0.9853} + 6\tabcolsep}@{}}{%
POSTGRES\_PORT \textbar{} int \textbar{} 5432 \textbar{} Puerto de PostgreSQL \textbar{}} \\
\multicolumn{4}{@{}>{\raggedright\arraybackslash}p{(\linewidth - 6\tabcolsep) * \real{0.9853} + 6\tabcolsep}@{}}{%
POSTGRES\_PASS \textbar{} str \textbar{} ``\,'' \textbar{} Contraseña de PostgreSQL \textbar{}} \\
\multicolumn{4}{@{}>{\raggedright\arraybackslash}p{(\linewidth - 6\tabcolsep) * \real{0.9853} + 6\tabcolsep}@{}}{%
ACCESS\_TOKEN \textbar{} int \textbar{} 11520 \textbar{} Expiración del JWT (8 días) \textbar{}} \\
\multicolumn{4}{@{}>{\raggedright\arraybackslash}p{(\linewidth - 6\tabcolsep) * \real{0.9853} + 6\tabcolsep}@{}}{%
FRONTEND\_HOST \textbar{} str \textbar{} http://localhost:5173 \textbar{} URL del frontend para CORS \textbar{}} \\
CORS\_ORIGINS \textbar{} list{[}str{]} \textbar{} {[}{]} & \multicolumn{3}{>{\raggedright\arraybackslash}p{(\linewidth - 6\tabcolsep) * \real{0.7500} + 4\tabcolsep}@{}}{%
Orígenes adicionales CORS \textbar{}} \\
\multicolumn{4}{@{}>{\raggedright\arraybackslash}p{(\linewidth - 6\tabcolsep) * \real{0.9853} + 6\tabcolsep}@{}}{%
ENVIRONMENT \textbar{} Literal \textbar{} ``local'' \textbar{} local \textbar{} staging \textbar{} production \textbar{}} \\
\multicolumn{4}{@{}>{\raggedright\arraybackslash}p{(\linewidth - 6\tabcolsep) * \real{0.9853} + 6\tabcolsep}@{}}{%
SMTP\_HOST \textbar{} str \textbar{} None \textbar{} None \textbar{} Host SMTP para envío de emails \textbar{}} \\
\multicolumn{4}{@{}>{\raggedright\arraybackslash}p{(\linewidth - 6\tabcolsep) * \real{0.9853} + 6\tabcolsep}@{}}{%
SMTP\_PORT \textbar{} int \textbar{} 587 \textbar{} Puerto SMTP \textbar{}} \\
\multicolumn{4}{@{}>{\raggedright\arraybackslash}p{(\linewidth - 6\tabcolsep) * \real{0.9853} + 6\tabcolsep}@{}}{%
SMTP\_TLS \textbar{} bool \textbar{} True \textbar{} Usar TLS en SMTP \textbar{}} \\
\multicolumn{4}{@{}>{\raggedright\arraybackslash}p{(\linewidth - 6\tabcolsep) * \real{0.9853} + 6\tabcolsep}@{}}{%
SMTP\_SSL \textbar{} bool \textbar{} False \textbar{} Usar SSL en SMTP \textbar{}} \\
\multicolumn{4}{@{}>{\raggedright\arraybackslash}p{(\linewidth - 6\tabcolsep) * \real{0.9853} + 6\tabcolsep}@{}}{%
SMTP\_USER \textbar{} str \textbar{} None \textbar{} None \textbar{} Usuario SMTP \textbar{}} \\
\multicolumn{4}{@{}>{\raggedright\arraybackslash}p{(\linewidth - 6\tabcolsep) * \real{0.9853} + 6\tabcolsep}@{}}{%
SMTP\_PASSWORD \textbar{} str \textbar{} None \textbar{} None \textbar{} Contraseña SMTP \textbar{}} \\
\multicolumn{4}{@{}>{\raggedright\arraybackslash}p{(\linewidth - 6\tabcolsep) * \real{0.9853} + 6\tabcolsep}@{}}{%
EMAILS\_FROM \textbar{} EmailStr \textbar{} None \textbar{} None \textbar{} Email remitente \textbar{}} \\
\multicolumn{4}{@{}>{\raggedright\arraybackslash}p{(\linewidth - 6\tabcolsep) * \real{0.9853} + 6\tabcolsep}@{}}{%
EMAIL\_RESET\_TOKEN \textbar{} int \textbar{} 48 \textbar{} Expiración del token de reset \textbar{}} \\
\multicolumn{4}{@{}>{\raggedright\arraybackslash}p{(\linewidth - 6\tabcolsep) * \real{0.9853} + 6\tabcolsep}@{}}{%
SENTRY\_DSN \textbar{} HttpUrl \textbar{} None \textbar{} None \textbar{} DSN de Sentry para error tracking \textbar{}} \\
\end{longtable}
}

\subsection{\texorpdfstring{Propiedades calculadas (no en \texttt{.env})}{Propiedades calculadas (no en .env)}}\label{propiedades-calculadas-no-en-.env}

{\def\LTcaptype{none} % do not increment counter
\begin{longtable}[]{@{}
  >{\raggedright\arraybackslash}p{(\linewidth - 2\tabcolsep) * \real{0.2874}}
  >{\raggedright\arraybackslash}p{(\linewidth - 2\tabcolsep) * \real{0.7126}}@{}}
\toprule\noalign{}
\begin{minipage}[b]{\linewidth}\raggedright
Propiedad
\end{minipage} & \begin{minipage}[b]{\linewidth}\raggedright
Descripción
\end{minipage} \\
\midrule\noalign{}
\endhead
\bottomrule\noalign{}
\endlastfoot
\texttt{SQLALCHEMY\_DATABASE\_URI} & URI completa de conexión a PostgreSQL \\
\texttt{all\_cors\_origins} & \texttt{BACKEND\_CORS\_ORIGINS} + \texttt{FRONTEND\_HOST} \\
\texttt{emails\_enabled} & \texttt{True} si \texttt{SMTP\_HOST} y \texttt{EMAILS\_FROM\_EMAIL} están configurados \\
\end{longtable}
}

\section{Convenciones de Código}\label{convenciones-de-cuxf3digo}

\subsection{Backend (Python)}\label{backend-python}

\textbf{Orden de imports:}

\begin{Shaded}
\begin{Highlighting}[]
\CommentTok{\# 1. Biblioteca estándar}
\ImportTok{import}\NormalTok{ uuid}
\ImportTok{from}\NormalTok{ typing }\ImportTok{import}\NormalTok{ Any}
\ImportTok{from}\NormalTok{ decimal }\ImportTok{import}\NormalTok{ Decimal}

\CommentTok{\# 2. Terceros}
\ImportTok{from}\NormalTok{ fastapi }\ImportTok{import}\NormalTok{ APIRouter, HTTPException}
\ImportTok{from}\NormalTok{ sqlmodel }\ImportTok{import}\NormalTok{ SQLModel, Field}

\CommentTok{\# 3. Locales}
\ImportTok{from}\NormalTok{ app.api.deps }\ImportTok{import}\NormalTok{ CurrentUser, SessionDep}
\ImportTok{from}\NormalTok{ app.models }\ImportTok{import}\NormalTok{ ProductCreate, ProductPublic}
\end{Highlighting}
\end{Shaded}

\textbf{Nomenclatura:}

\begin{itemize}
\tightlist
\item
  Funciones y variables: \texttt{snake\_case}
\item
  Clases y modelos: \texttt{PascalCase}
\item
  Modelos de BD: singular (\texttt{Product}, \texttt{Sale})
\item
  Schemas: \texttt{\{Modelo\}Create}, \texttt{\{Modelo\}Update}, \texttt{\{Modelo\}Public}, \texttt{\{Modelo\}sPublic}
\item
  Aliases de tipo: \texttt{PascalCase} con \texttt{Annotated} (\texttt{CurrentUser}, \texttt{SessionDep})
\end{itemize}

\textbf{Anotaciones de tipo:} obligatorias en todas las funciones. Usar \texttt{str\ \textbar{}\ None} (no \texttt{Optional{[}str{]}}).

\textbf{Patrón de rutas:}

\begin{Shaded}
\begin{Highlighting}[]
\AttributeTok{@router.post}\NormalTok{(}\StringTok{"/"}\NormalTok{, response\_model}\OperatorTok{=}\NormalTok{ProductPublic)}
\KeywordTok{def}\NormalTok{ create\_product(}
    \OperatorTok{*}\NormalTok{,}
\NormalTok{    session: SessionDep,}
\NormalTok{    current\_user: CurrentUser,}
\NormalTok{    current\_organization: CurrentOrganization,}
\NormalTok{    product\_in: ProductCreate,}
\NormalTok{) }\OperatorTok{{-}\textgreater{}}\NormalTok{ Any:}
\NormalTok{    product }\OperatorTok{=}\NormalTok{ Product.model\_validate(product\_in, update}\OperatorTok{=}\NormalTok{\{}\StringTok{"organization\_id"}\NormalTok{: current\_organization\})}
\NormalTok{    session.add(product)}
\NormalTok{    session.commit()}
\NormalTok{    session.refresh(product)}
    \ControlFlowTok{return}\NormalTok{ product}
\end{Highlighting}
\end{Shaded}

\textbf{Códigos HTTP usados:}

\begin{itemize}
\tightlist
\item
  \texttt{400}: Solicitud inválida (datos incorrectos)
\item
  \texttt{403}: Sin permisos / credenciales inválidas
\item
  \texttt{404}: Recurso no encontrado
\item
  \texttt{409}: Conflicto (recurso duplicado)
\end{itemize}

\textbf{Multi-tenancy:} \textbf{todas} las consultas de negocio incluyen filtro \texttt{organization\_id\ =\ current\_organization}. El valor se extrae del usuario autenticado, nunca del cuerpo del request.

\subsection{Frontend (TypeScript/React)}\label{frontend-typescriptreact}

\textbf{Orden de imports:}

\begin{Shaded}
\begin{Highlighting}[]
\CommentTok{// 1. Externos}
\ImportTok{import}\NormalTok{ \{ useState \} }\ImportTok{from} \StringTok{"react"}\OperatorTok{;}
\ImportTok{import}\NormalTok{ \{ useMutation}\OperatorTok{,}\NormalTok{ useQueryClient \} }\ImportTok{from} \StringTok{"@tanstack/react{-}query"}\OperatorTok{;}
\ImportTok{import}\NormalTok{ \{ z \} }\ImportTok{from} \StringTok{"zod"}\OperatorTok{;}

\CommentTok{// 2. Internos (alias @/)}
\ImportTok{import}\NormalTok{ \{ ProductsService \} }\ImportTok{from} \StringTok{"@/client"}\OperatorTok{;}
\ImportTok{import}\NormalTok{ \{ DataTable \} }\ImportTok{from} \StringTok{"@/components/Common/DataTable"}\OperatorTok{;}
\ImportTok{import}\NormalTok{ \{ useCustomToast \} }\ImportTok{from} \StringTok{"@/hooks/useCustomToast"}\OperatorTok{;}
\end{Highlighting}
\end{Shaded}

\textbf{Componentes:} solo funcionales, un componente por archivo, nombre de archivo en \texttt{PascalCase}.

\textbf{Patrón de mutación:}

\begin{Shaded}
\begin{Highlighting}[]
\KeywordTok{const}\NormalTok{ mutation }\OperatorTok{=} \FunctionTok{useMutation}\NormalTok{(\{}
\NormalTok{  mutationFn}\OperatorTok{:}\NormalTok{ (data}\OperatorTok{:}\NormalTok{ ProductCreate) }\KeywordTok{=\textgreater{}}
\NormalTok{    ProductsService}\OperatorTok{.}\FunctionTok{createProduct}\NormalTok{(\{ requestBody}\OperatorTok{:}\NormalTok{ data \})}\OperatorTok{,}
\NormalTok{  onSuccess}\OperatorTok{:}\NormalTok{ () }\KeywordTok{=\textgreater{}} \FunctionTok{showSuccessToast}\NormalTok{(}\StringTok{"Producto creado"}\NormalTok{)}\OperatorTok{,}
\NormalTok{  onError}\OperatorTok{:}\NormalTok{ (err) }\KeywordTok{=\textgreater{}} \FunctionTok{handleError}\NormalTok{(err)}\OperatorTok{,}
\NormalTok{  onSettled}\OperatorTok{:}\NormalTok{ () }\KeywordTok{=\textgreater{}}\NormalTok{ queryClient}\OperatorTok{.}\FunctionTok{invalidateQueries}\NormalTok{(\{ queryKey}\OperatorTok{:}\NormalTok{ [}\StringTok{"products"}\NormalTok{] \})}\OperatorTok{,}
\NormalTok{\})}\OperatorTok{;}
\end{Highlighting}
\end{Shaded}

\textbf{Validación de formularios:} React Hook Form + Zod schemas con \texttt{zodResolver}.

\textbf{No editar:}

\begin{itemize}
\tightlist
\item
  \texttt{src/client/**}: auto-generado desde OpenAPI
\item
  \texttt{src/components/ui/**}: componentes Shadcn/ui
\item
  \texttt{src/routeTree.gen.ts}: auto-generado por TanStack Router
\end{itemize}

\section{Pipeline de CI/CD}\label{pipeline-de-cicd}

El repositorio incluye workflows de GitHub Actions:

\subsection{\texorpdfstring{Workflow de pruebas (en cada PR y push a \texttt{main})}{Workflow de pruebas (en cada PR y push a main)}}\label{workflow-de-pruebas-en-cada-pr-y-push-a-main}

\begin{enumerate}
\def\labelenumi{\arabic{enumi}.}
\tightlist
\item
  \textbf{Backend:} Ejecuta \texttt{uv\ run\ bash\ scripts/test.sh}; incluye Mypy, Ruff y Pytest con cobertura.
\item
  \textbf{Frontend:} Ejecuta \texttt{bun\ run\ lint}; Biome linting y verificación de tipos TypeScript.
\end{enumerate}

\subsection{Workflow de despliegue}\label{workflow-de-despliegue}

Al hacer merge a \texttt{main}:

\begin{enumerate}
\def\labelenumi{\arabic{enumi}.}
\tightlist
\item
  Construye las imágenes Docker de backend y frontend.
\item
  Las etiqueta con el hash del commit y con \texttt{latest}.
\item
  Las publica en el registro de imágenes.
\item
  Actualiza los contenedores en el servidor de producción via SSH.
\end{enumerate}

\section{Estructura de Pruebas}\label{estructura-de-pruebas}

\subsection{Backend (pytest)}\label{backend-pytest}

\begin{verbatim}
backend/tests/
├── conftest.py                     # Fixtures globales: db, client, auth headers
├── api/routes/
│   ├── test_login.py               # Autenticación, recuperación de contraseña
│   ├── test_users.py               # CRUD de usuarios, cambio de contraseña
│   ├── test_categories.py          # CRUD de categorías
│   ├── test_products.py            # CRUD de productos, ajuste de stock
│   ├── test_customers.py           # CRUD de clientes, historial de compras
│   ├── test_sales.py               # Registro y cancelación de ventas
│   ├── test_inventory_movements.py # Listado de movimientos
│   ├── test_dashboard.py           # Estadísticas del dashboard
│   └── test_private.py             # Endpoints de superusuario
└── crud/
    ├── test_user.py
    ├── test_category.py
    ├── test_product.py
    ├── test_customer.py
    ├── test_sale.py
    ├── test_inventory_movement.py
    └── test_dashboard.py
\end{verbatim}

\textbf{Fixtures disponibles en conftest.py:}

{\def\LTcaptype{none} % do not increment counter
\begin{longtable}[]{@{}
  >{\raggedright\arraybackslash}p{(\linewidth - 2\tabcolsep) * \real{0.3462}}
  >{\raggedright\arraybackslash}p{(\linewidth - 2\tabcolsep) * \real{0.6538}}@{}}
\toprule\noalign{}
\begin{minipage}[b]{\linewidth}\raggedright
Fixture
\end{minipage} & \begin{minipage}[b]{\linewidth}\raggedright
Descripción
\end{minipage} \\
\midrule\noalign{}
\endhead
\bottomrule\noalign{}
\endlastfoot
\texttt{db} & Sesión de BD para pruebas (con rollback automático) \\
\texttt{client} & \texttt{TestClient} de FastAPI \\
\texttt{superuser\_token\_headers} & Headers con token del superusuario \\
\texttt{normal\_user\_token\_headers} & Headers con token de usuario normal \\
\end{longtable}
}

\textbf{Ejecutar pruebas:}

\begin{Shaded}
\begin{Highlighting}[]
\BuiltInTok{cd}\NormalTok{ backend}
\ExtensionTok{uv}\NormalTok{ run bash scripts/test.sh           }\CommentTok{\# Todos los tests con cobertura}
\ExtensionTok{uv}\NormalTok{ run pytest tests/api/routes/ }\AttributeTok{{-}v}    \CommentTok{\# Solo tests de API}
\ExtensionTok{uv}\NormalTok{ run pytest }\AttributeTok{{-}k} \StringTok{"test\_create"} \AttributeTok{{-}v}     \CommentTok{\# Tests por patrón}
\end{Highlighting}
\end{Shaded}

\subsection{Frontend (Playwright E2E)}\label{frontend-playwright-e2e}

\begin{verbatim}
frontend/tests/
├── auth.setup.ts                   # Configuración de autenticación compartida
├── config.ts                       # Variables de configuración de tests
├── login.spec.ts                   # Flujo de inicio de sesión
├── sign-up.spec.ts                 # Registro de usuario
├── organization-signup.spec.ts     # Registro de organización
├── reset-password.spec.ts          # Recuperación de contraseña
├── inventory.spec.ts               # Módulo de inventario
├── sales.spec.ts                   # Módulo de ventas
├── customers.spec.ts               # Módulo de clientes
├── admin.spec.ts                   # Gestión de usuarios
├── user-settings.spec.ts           # Configuración de cuenta
└── utils/                          # Helpers de pruebas E2E
\end{verbatim}

\section{Generación del Cliente API}\label{generaciuxf3n-del-cliente-api}

El cliente TypeScript del frontend se genera automáticamente desde el esquema OpenAPI del backend:

\begin{Shaded}
\begin{Highlighting}[]
\CommentTok{\# El backend debe estar corriendo en http://localhost:8000}
\BuiltInTok{cd}\NormalTok{ frontend}
\ExtensionTok{bun}\NormalTok{ run generate{-}client}
\end{Highlighting}
\end{Shaded}

La configuración del generador está en \texttt{frontend/openapi-ts.config.ts}. Los archivos generados en \texttt{src/client/} no deben editarse manualmente; cualquier cambio se sobreescribirá en la próxima generación.

\section{Multi-Tenancy}\label{multi-tenancy}

OrbitEngine implementa multi-tenancy mediante \textbf{aislamiento por discriminador en tablas compartidas} (Shared Database, Separate Rows):

\begin{itemize}
\tightlist
\item
  Todas las tablas de negocio (\texttt{categories}, \texttt{products}, \texttt{customers}, \texttt{sales}, \texttt{sale\_items}, \texttt{inventory\_movements}) incluyen la columna \texttt{organization\_id}.
\item
  La \texttt{organization\_id} se indexa en cada tabla para garantizar rendimiento en consultas filtradas.
\item
  Toda consulta de datos de negocio filtra por \texttt{organization\_id\ =\ current\_user.organization\_id}.
\item
  El valor de \texttt{organization\_id} \textbf{nunca} se acepta como parámetro del request; siempre se extrae del token JWT del usuario autenticado.
\item
  Las restricciones de unicidad incluyen \texttt{organization\_id} como parte de la clave compuesta (ej. \texttt{UNIQUE(organization\_id,\ sku)}), permitiendo que el mismo SKU exista en organizaciones distintas.
\end{itemize}

Esta estrategia garantiza aislamiento completo de datos entre tenants sin necesidad de esquemas o bases de datos separadas, facilitando el mantenimiento y la escalabilidad horizontal.

\emph{Documento generado como parte del proyecto de grado. Universidad Sergio Arboleda, Semillero de Software como Innovación, Abril 2026.}
